\section{第一性原理：从 ``神秘 AGI'' 到 ``可分解系统''}

\subsection{输入法类比：复杂智能可由简单模块组合}
Yangqing Jia 在开场强调，外界常把 AGI 描述为神秘黑箱，但工程视角更强调可分解性。他用中文输入法做类比：用户输入拼音后，系统并不是 ``一次性神谕输出''，而是经历切分、候选召回、排序、语言模型修正等分层模块，再由交互反馈不断改进体验。

这个类比指向一个关键判断：大模型系统同样可以被拆为相对清晰的子问题。即使底层模型参数巨大，上层产品与基础设施仍然需要可解释的模块边界，否则无法调试、优化与治理。

\begin{importantbox}{工程含义：把 ``智能'' 还原为 ``流水线''}
输入法并不需要 ``理解一切'' 才能工作良好，它依赖高质量候选生成与排序。Agent 系统也类似，不必追求单模型包办所有功能，而应优先把召回、规划、执行、验证各阶段做到可控和可观测。
\end{importantbox}

\begin{knowledgebox}{输入法类比映射到 Agent Pipeline}
\begin{itemize}
    \item 拼音切分 $\leftrightarrow$ 任务解析（Task Parsing）；
    \item 候选字召回 $\leftrightarrow$ 工具/知识候选检索（Retrieval）；
    \item 候选排序 $\leftrightarrow$ 规划与动作选择（Planning / Policy）；
    \item 用户纠错 $\leftrightarrow$ 在线反馈与后训练闭环（Feedback Loop）。
\end{itemize}
\end{knowledgebox}

\begin{warningbox}{不要把 ``可分解'' 误解为 ``简单''}
可分解并不意味着成本低。相反，模块越多，跨模块协议、数据一致性与故障联动风险越高。工程团队必须在 ``模块化收益'' 与 ``系统复杂度'' 之间做严谨权衡。
\end{warningbox}

\subsection{本章小结}
``神秘能力'' 在工程上最终要落到 ``可分解系统''。这不是降低难度，而是把难点从模型单点，转移到跨模块协同与系统可靠性。

